\BourbakiStarChapter{To the Reader}

\begin{enumerate}
\def\labelenumi{\arabic{enumi}.}
\item
  This series of volumes, a list of which is given on pages ix and x, takes up mathematics at the beginning, and gives complete proofs. In principle, it requires no particular knowledge of mathematics on the readers' part, but only a certain familiarity with mathematical reasoning and a certain capacity for abstract thought. Nevertheless, it is directed especially to those who have a good knowledge of at least the content of the first year or two of a university mathematics course.
\item
  The method of exposition we have chosen is axiomatic and abstract, and normally proceeds from the general to the particular. This choice has been dictated by the main purpose of the treatise, which is to provide a solid foundation for the whole body of modern mathematics. For this it is indispensable to become familiar with a rather large number of very general ideas and principles. Moreover, the demands of proof impose a rigorously fixed order on the subject matter. It follows that the utility of certain considerations will not be immediately apparent to the reader unless he has already a fairly extended knowledge of mathematics; otherwise he must have the patience to suspend judgment until the occasion arises.
\item
  In order to mitigate this disadvantage we have frequently inserted examples in the text which refer to facts the reader may already know but which have not yet been discussed in the series. Such examples are always placed between two asterisks: *\ldots*. Most readers will undoubtedly find that these examples will help them to understand the text, and will prefer not to leave them out, even at a first reading. Their omission would of course have no disadvantage, from a purely logical point of view.
\item
  This series is divided into volumes (here called ``Books''). The first six Books are numbered and, in general, every statement in the text assumes as known only those results which have already been discussed in the preceding volumes. This rule holds good within each Book, but for convenience of exposition these Books are no longer arranged in a consecutive order. At the beginning of each of these Books (or of these chapters), the reader will find a precise indication of its logical relationship to the other Books and he will thus be able to satisfy himself of the absence of any vicious circle.
\item
  The logical framework of each chapter consists of the definitions, the axioms, and the theorems of the chapter. These are the parts that have mainly to be borne in mind for subsequent use. Less important results and those which can easily be deduced from the theorems are labelled as ``propositions'', ``lemmas'', ``corollaries'', ``remarks'', etc. Those which may be omitted at a first reading are printed in small type. A commentary on a particularly important theorem appears occasionally under the name of ``scholium''.
\end{enumerate}

To avoid tedious repetitions it is sometimes convenient to introduce notations or abbreviations which are in force only within a certain chapter or a certain section of a chapter (for example, in a chapter which is concerned only with commutative rings, the word ``ring'' would always signify ``commutative ring''). Such conventions are always explicitly mentioned, generally at the beginning of the chapter in which they occur.

\begin{enumerate}
\def\labelenumi{\arabic{enumi}.}
\setcounter{enumi}{5}
\item
  Some passages in the text are designed to forewarn the reader against serious errors. These passages are signposted in the margin with the sign Z (``dangerous bend'').
\item
  The Exercises are designed both to enable the reader to satisfy himself that he has digested the text and to bring to his notice results which have no place in the text but which are nonetheless of interest. The most difficult exercises bear the sign ¶.
\item
  In general, we have adhered to the commonly accepted terminology, except where there appeared to be good reasons for deviating from it.
\item
  We have made a particular effort always to use rigorously correct language, without sacrificing simplicity. As far as possible we have drawn attention in the text to abuses of language, without which any mathematical text runs the risk of pedantry, not to say unreadability.
\item
  Since in principle the text consists of the dogmatic exposition of a theory, it contains in general no references to the literature. Bibliographical references are gathered together in Historical Notes, usually at the end of each chapter. These notes also contain indications, where appropriate, of the unsolved problems of the theory.
\end{enumerate}

The bibliography which follows each historical note contains in general only those books and original memoirs which have been of the greatest importance in the evolution of the theory under discussion. It makes no sort of pretence to completeness; in particular, references which serve only to determine questions of priority are almost always omitted.

As to the exercises, we have not thought it worthwhile in general to indicate their origins, since they have been taken from many different sources (original papers, textbooks, collections of exercises).

\par\noindent\textbf{11. References to a part of this series are given as follows:}\par

\begin{enumerate}
\def\labelenumi{\alph{enumi})}
\item
  If reference is made to theorems, axioms, or definitions presented in the same section, they are quoted by their number.
\item
  If they occur in another section of the same chapter, this section is also quoted in the reference.
\item
  If they occur in another chapter in the same Book, the chapter and section are quoted.
\item
  If they occur in another Book, this Book is first quoted by its title.
\end{enumerate}

The Summaries of Results are quoted by the letter R; thus Set Theory, R signifies ``Summary of Results of the Theory of Sets''.

\BourbakiStarChapter{Contents of the Elements of Mathematics Series}

\BourbakiSeriesBook{Theory of Sets}

\begin{enumerate}
\def\labelenumi{\arabic{enumi}.}
\tightlist
\item
  Description of formal mathematics. 2. Theory of sets. 3. Ordered sets; cardinals; natural numbers. 4. Structures.
\end{enumerate}

\BourbakiSeriesBook{Algebra}

\begin{enumerate}
\def\labelenumi{\arabic{enumi}.}
\tightlist
\item
  Algebraic structures. 2. Linear algebra. 3. Tensor algebras, exterior algebras, symmetric algebras. 4. Polynomials and rational fractions. 5. Fields. 6. Ordered groups and fields. 7. Modules over principal ideal rings. 8. Semi-simple modules and rings. 9. Sesquilinear and quadratic forms.
\end{enumerate}

\BourbakiSeriesBook{General Topology}

\begin{enumerate}
\def\labelenumi{\arabic{enumi}.}
\tightlist
\item
  Topological structures. 2. Uniform structures. 3. Topological groups. 4. Real numbers. 5. One-parameter groups. 6. Real number spaces, affine and projective spaces. 7. The additive groups \(\mathbf{R}^n\) . 8. Complex numbers. 9. Use of real numbers in general topology. 10. Function spaces.
\end{enumerate}

\BourbakiSeriesBook{Functions of a Real Variable}

\begin{enumerate}
\def\labelenumi{\arabic{enumi}.}
\tightlist
\item
  Derivatives. 2. Primitives and integrals. 3. Elementary functions. 4. Differential equations. 5. Local study of functions. 6. Generalized Taylor expansions. The Euler-Maclaurin summation formula. 7. The gamma function. Dictionary.
\end{enumerate}

\BourbakiSeriesBook{Topological Vector Spaces}

\begin{enumerate}
\def\labelenumi{\arabic{enumi}.}
\tightlist
\item
  Topological vector spaces over a valued field. 2. Convex sets and locally convex spaces. 3. Spaces of continuous linear mappings. 4. Duality in topological vector spaces. 5. Hilbert spaces: elementary theory. Dictionary.
\end{enumerate}

\BourbakiSeriesBook{Integration}

\begin{enumerate}
\def\labelenumi{\arabic{enumi}.}
\tightlist
\item
  Convexity inequalities. 2. Riesz spaces. 3. Measures on locally compact spaces. 4. Extension of a measure. L \(^{p}\) spaces. 5. Integration of measures. 6. Vectorial integration. 7. Haar measure. 8. Convolution and representation.
\end{enumerate}

\BourbakiSeriesTitle{Lie Groups and Lie Algebras}

\begin{enumerate}
\def\labelenumi{\arabic{enumi}.}
\tightlist
\item
  Lie algebras. 2. Free Lie algebras. 3. Lie groups. 4. Coxeter groups and Tits systems. 5. Groups generated by reflections. 6. Root systems.
\end{enumerate}

\BourbakiSeriesTitle{Commutative Algebra}

\begin{enumerate}
\def\labelenumi{\arabic{enumi}.}
\tightlist
\item
  Flat modules. 2. Localization. 3. Graduations, filtrations, and topologies. 4. Associated prime ideals and primary decomposition. 5. Integers. 6. Valuations. 7. Divisors.
\end{enumerate}

\BourbakiSeriesTitle{Spectral Theory}

\begin{enumerate}
\def\labelenumi{\arabic{enumi}.}
\tightlist
\item
  Normed algebras. 2. Locally compact groups.
\end{enumerate}

\BourbakiSeriesTitle{Differential and Analytic Manifolds}

Summary of results.

\BourbakiStarChapter{Introduction}

Algebra is essentially concerned with calculating, that is, performing, on elements of a set, ``algebraic operations'', the most well-known example of which is provided by the ``four rules'' of elementary arithmetic.

This is not the place to retrace the slow process of progressive abstraction by which the notion of algebraic operation, at first restricted to natural numbers and measurable magnitudes, little by little enlarged its domain, whilst the notion of ``number'' received a parallel generalization, until, passing beyond the latter, it came to be applied to elements with no ``numerical'' character, for example permutations of a set (see the Historical Note to Chapter I). It is no doubt the possibility of these successive extensions, in which the form of the calculations remained the same, whereas the nature of the mathematical entities subjected to these calculations varied considerably, which was responsible for the gradual isolation of the guiding principle of modern mathematics, namely that mathematical entities in themselves are of little importance; what matters are their relations (see Book I). It is certain in any case that Algebra attained this level of abstraction well before the other branches of Mathematics and it has long been accustomed to being considered as the study of algebraic operations, independent of the mathematical entities to which they can be applied.

Deprived of any specific character, the common notion underlying the usual algebraic operations is very simple: performing an algebraic operation on two elements \(a, b\) of the same set \(\mathbf{E}\) , means associating with the ordered pair \((a, b)\) a well-defined third element \(c\) of the set \(\mathbf{E}\) . In other words, there is nothing more in this notion than that of function: to be given an algebraic operation is to be given a function defined on \(\mathbf{E} \times \mathbf{E}\) and taking its values in \(\mathbf{E}\) ; the only peculiarity lies in the fact that the set of definition of the function is the product of two sets identical with the set where the function takes its values; to such a function we give the name law of composition.

Alongside these ``internal laws'', mathematicians have been led (chiefly under the influence of Geometry) to consider another type of ``law of composition''; these are the ``laws of action'', where, besides the set E (which remains as it were in the foreground), an auxiliary set \(\Omega\) occurs, whose elements are called operators: the law this time associating with an ordered pair \((\alpha, a)\) consisting of an operator \(\alpha \in \Omega\) and an element \(a \in E\) a second element b of E. For example a homothety of given centre on the Euclidean space E associates with a real number k (the ``homothety ratio'', which is here the operator) and a point A of E another point A' of E: this is a law of action on E.

In conformity with the general definitions (Set Theory, IV, § 1, no. 4), being given on a set E one or several laws of composition or laws of action defines a structure on E; for the structures defined in this way we reserve precisely the name algebraic structures and it is the study of these which constitutes Algebra.

There are several species (Set Theory, IV, § 1, no. 4) of algebraic structures, characterized on the one hand by the laws of composition or laws of action which define them and on the other by the axioms to which these laws are subjected. Of course, these axioms have not been chosen arbitrarily, but are just the properties of most of the laws which occur in applications, such as associativity, commutativity, etc. Chapter I is essentially devoted to the exposition of these axioms and the general consequences which follow from them; also there is a more detailed study of the two most important species of algebraic structure, that of a group (in which only one law of composition occurs) and that of a ring (with two laws of composition), of which a field structure is a special case.

In Chapter I are also found the definitions of groups with operators and rings with operators, where, besides the laws of composition, there occur one or several laws of action. The most important groups with operators are modules, of which vector spaces are a particular example, which play an important role both in Classical Geometry and in Modern Analysis. The study of module structures derives its origin from that of linear equations, whence its name Linear Algebra; the results on this subject are to be found in Chapter II.

Similarly, the rings with operators which occur most often are called algebras (or hypercomplex systems). In Chapters III and IV a detailed study is made of two particular types of algebras: exterior algebras, which with the theory of determinants which is derived from them is a valuable tool of Linear Algebra; and polynomial algebras, which are fundamental to the theory of algebraic equations.

In Chapter V there is an exposition of the general theory of commutative fields and their classification. The origin of this theory is the study of algebraic equations in one unknown; the questions which gave birth to this today have little more than a historical interest, but the theory of commutative fields remains fundamental to Algebra, being basic to the theory of algebraic numbers on the one hand and Algebraic Geometry on the other.

As the set of natural numbers has two laws of composition, addition and multiplication, Classical Arithmetic (or Number Theory), which is the study of natural numbers, is subordinate to Algebra. However, related to the algebraic structure defined by these two laws is the structure defined by the order relation ``a divides b''; and the distinguishing aspect of Classical Arithmetic is precisely the study of the relations between these two associated structures. This is not the only example where an order structure is thus associated with an algebraic structure by a ``divisibility'' relation: this relation plays just as important a role in polynomial rings. A general study will therefore be made of this in Chapter VI; this study will be applied in Chapter VII to determining the structure of modules over certain particularly simple rings and in particular to the theory of ``elementary divisors''.

Chapters VIII and IX are devoted to more particular theories, but which have many applications in Analysis: on the one hand, the theory of semisimple modules and rings, closely related to that of linear representations of groups; on the other, the theory of quadratic forms and Hermitian forms, with the study of the groups associated with them. Finally, the elementary geometries (affine, projective, euclidean, etc.) are studied in Chapters II, VI and IX from a purely algebraic point of view: here there is little more than a new language for expressing results of Algebra already obtained elsewhere, but it is a language which is particularly well adapted to the later developments of Algebraic Geometry and Differential Geometry, to which this chapter serves as an introduction.
